\documentclass{article}
\usepackage[UTF8]{ctex}
\usepackage{amsmath,amssymb,amsfonts}
\usepackage{enumitem}
\usepackage{geometry}
\geometry{a4paper, margin=1in}
% === 错题本排版优化（不参与题目标识） ===
\setlength{\parindent}{0pt}
\setlength{\parskip}{0.35em plus 0.10em minus 0.08em}
\setlength{\emergencystretch}{3em}
\setlist{itemsep=0.25em, topsep=0.40em, parsep=0pt}
\allowdisplaybreaks[2]
% === 错题本排版优化结束 ===
\begin{document}
\section*{郑州大学2024年数学分析真题}
\section*{127．郑州大学2024年数学分析试卷}

\subsection{题目1}
一．（10 分）求函数极限 \(\lim _{x \rightarrow 0^{+}}\left[\frac{1}{x(1+x)^{\frac{1}{x}}}-\frac{1}{x e}\right]\) ．

\subsection{题目2}
二．（10 分）求定积分 \(\int_{0}^{2024} \frac{x}{e^{2024-x}+c^{x}} \mathrm{~d} x\) ．

\subsection{题目3}
三．（10 分）设椭圆 \(\frac{x^{2}}{4}+\frac{y^{2}}{9}=1\) 在点 \(A\left(1, \frac{3 \sqrt{3}}{2}\right)\) 的切线交 \(y\) 轴于点 \(B, L\) 是从 \(A\) 到 \(B\) 的直线段，计算

\[
\int_{L}\left(\frac{\sin y}{x+1}-\sqrt{3} y\right) \mathrm{d} x+(\cos y \ln (1+x)+2 \sqrt{3} x-\sqrt{3}) \mathrm{d} y .
\]

\subsection{题目4}
四．（10分）讨论 \(\int_{0}^{+\infty} \frac{\sin x}{x^{p-1}+\frac{1}{x}} \mathrm{~d} x(p \geq 0)\) 的条件收敛和绝对收敛性．

\subsection{题目5}
五．（10 分）求函数 \(f(x)=\int_{0}^{x} \frac{\ln (1+2 t)}{t} \mathrm{~d} t\) 的麦克劳林级数展开式．

\subsection{题目6}
六．（15 分）讨论方程 \(f(x)=-\frac{1}{2}\left(1+e^{-1}\right)+\int_{-1}^{1}|x-t| e^{-t^{2}} \mathrm{~d} t=0\) 在［－1，1］上根的个数．

\subsection{题目7}
七．（15 分）设 \(f(x, y)=\left\{\begin{array}{ll}\frac{|x|^{\alpha}|y|^{\alpha}}{x^{2}+y^{2}}, & x^{2}+y^{2} \neq 0 ; \\ 0, & x^{2}+y^{2}=0 .\end{array}\right.\) 证明：

\begin{itemize}
  \item[（1）] \(\alpha>1\) 时，\(f(x, y)\) 在 \((0,0)\) 连续．
  \item[（2）] \(\alpha>\frac{3}{2}\) 时，\(f(x, y)\) 在 \((0,0)\) 可微．
\end{itemize}

\subsection{题目8}
八．（15 分）证明：函数项级数 \(\sum_{n=1}^{\infty} \frac{x^{2}}{\left(1+x^{2}\right)^{n}}\) 在 \((-\infty,+\infty)\) 上点点收敛，但并非一致收敛。

\subsection{题目9}
九．（ 35 分）证明题．

\begin{itemize}
  \item[（1）] （5 分）已知 \(\lim _{x \rightarrow+\infty}[f(x)-(a x+b)]=0\) ，求 \(a, b\) ．
  \item[（2）] （10 分）设 \(f(x)\) 在 \([a,+\infty)\) 上一致连续，\(g(x) \in C([a,+\infty))\) ，且 \(\lim _{x \rightarrow+\infty}[f(x)-g(x)]=0\) ．证明： \(g(x)\) 在 \([a,+\infty)\) 上一致连续．
  \item[（3）] （5 分）用（2）的结论说明 \(f(x)=\frac{1}{x}+\ln \left(1+e^{x}\right)\) 在 \([1,+\infty)\) 上一致连续．
  \item[（4）] （15 分）设函数 \(f(x)\) 在 \((-\infty,+\infty)\) 上可导，且存在常数 \(a, b, c, d(a<c)\) 使得
\[
\lim _{x \rightarrow-\infty}[f(x)-(a x+b)]=0, \lim _{x \rightarrow+\infty}[f(x)-(c x+d)]=0 .
\]
\end{itemize}

证明：对任意的 \(\lambda \in(a, c)\) ，存在 \(\xi \in(-\infty,+\infty)\) ，使得 \(f^{\prime}(\xi)=\lambda\) ．

\subsection{题目1101}
原题附加题：（20 分）证明：

\begin{itemize}
  \item[（1）] \(\int_{0}^{+\infty} \frac{\sin u x}{x} \mathrm{~d} x\) 在任何不包含 \(u=0\) 的闭区间 \([a, b]\) 上一致收敛。
  \item[（2）] 已知 \(\int_{0}^{+\infty} \frac{\sin x}{x} \mathrm{~d} x=\frac{\pi}{2}\) ，求 \(F(x)=\int_{0}^{+\infty} \frac{\sin ^{2} x t}{t^{2}} \mathrm{~d} t(x>0)\) ．
\end{itemize}
\end{document}
